
% ==================================================
\chapter{Introduction}
% ==================================================

\section{Overview}

LearnDeep is an educational deep learning framework
developed from scratch to explore the fundamental
concepts, algorithms, and engineering principles
behind modern deep learning systems.

The primary objective of this project is to understand
how a deep learning framework operates internally,
from tensor computations and numerical kernels to
automatic differentiation, neural network layers,
and model training.

Rather than treating a deep learning framework as a
black box, LearnDeep aims to make its underlying
mechanisms explicit through a modular implementation
and a focus on mathematical correctness.

\section{Motivation}

Modern deep learning frameworks provide high-level
interfaces for constructing, training, and evaluating
neural networks. However, many of the mechanisms
responsible for these operations remain hidden behind
their public APIs.

Developing LearnDeep provides an opportunity to
investigate these mechanisms directly, including:

\begin{itemize}
    \item Tensor representation, indexing, and
          memory ownership.
    \item Numerical operations, broadcasting, and
          CPU and GPU kernel execution.
    \item Computational graphs and reverse-mode
          automatic differentiation.
    \item Neural network layers, loss functions,
          and gradient propagation.
    \item Model training, optimization, and execution
          management.
\end{itemize}

Implementing these components from the ground up
provides practical insight into the relationship
between mathematical operations, software
architecture, and hardware-level computation.

\section{Scope}

LearnDeep is intended primarily as a learning and
experimentation platform. Its implementation
prioritizes understanding, modularity, and correctness
while exploring the design of a deep learning framework.

The technical reference documents the framework's
architecture, public interfaces, mathematical
foundations, and important implementation decisions.

It is organized around the major components of the
system, allowing readers to understand both how
individual components work and how they interact
during model execution and training.

\section{Organization of This Document}

The remainder of this document is organized as follows:

\begin{itemize}
    \item \textbf{Tensor System:} Tensor representation,
          shape semantics, data layout, and ownership.
    \item \textbf{Operations and Kernels:} Numerical
          operations, broadcasting, and backend execution.
    \item \textbf{Automatic Differentiation:}
          Computational graphs, gradient propagation,
          and backward execution.
    \item \textbf{Neural Network Components:} Layers,
          loss functions, and model construction.
    \item \textbf{Training:} Forward execution,
          optimization, and training workflows.
\end{itemize}

These sections describe the design and behavior of
the implemented system, including relevant
mathematical definitions and validation methods.
